\documentclass[aspectratio=169]{beamer}
\input{header}
\coursecode{JKSSB}

\title{RAM, ROM \& Cache}
\subtitle{Lesson 08 --- Volatile and Non-Volatile Memory}
\author{JKSSB Computer Awareness}
\date{\today}

\begin{document}
\frame{\titlepage}

\begin{frame}{Two Kinds of Memory}
  \begin{itemize}
    \item Main memory holds the data and instructions the CPU uses directly.
    \item The key split is \key{volatile} versus \key{non-volatile}.
    \item \key{Volatile}: contents are lost when power is removed.
    \item \key{Non-volatile}: contents are retained without power.
  \end{itemize}
  \begin{alertblock}{Trap (TRAP-01)}
    \bad{RAM is volatile; ROM is non-volatile.} Do not reverse the two.
  \end{alertblock}
\end{frame}

\begin{frame}{RAM --- Volatile Main Memory}
  \begin{itemize}
    \item \key{RAM} is \key{Random Access Memory}.
    \item It is the \key{main memory} that holds running programs and data.
    \item \key{Volatile}: on power-off, its contents disappear.
    \item ``Random access'' means any location is reached in roughly equal
      time.
    \item In the office desktop (EX-01), the OS, MS Word, and the spreadsheet
      load into RAM while in use.
  \end{itemize}
  \begin{block}{Exam cue}
    RAM = working space; cleared when the machine is switched off.
  \end{block}
\end{frame}

\begin{frame}{Two Types of RAM: DRAM and SRAM}
  \begin{columns}[T]
    \column{0.46\textwidth}
      \textbf{DRAM}
      \begin{itemize}
        \item Dynamic RAM
        \item Needs periodic \key{refresh}
        \item Denser and cheaper
        \item Used for \key{main memory}
      \end{itemize}
    \column{0.46\textwidth}
      \textbf{SRAM}
      \begin{itemize}
        \item Static RAM
        \item No refresh needed
        \item \key{Faster}, but costlier
        \item Used for \key{cache}
      \end{itemize}
  \end{columns}
  \vspace{0.2cm}
  \begin{alertblock}{Do not confuse}
    Both are RAM and both are volatile --- they differ in speed, cost, and role.
  \end{alertblock}
\end{frame}

\begin{frame}{DRAM vs SRAM at a Glance}
  \begin{center}
    \begin{tabular}{lll}
      \toprule
      \textbf{Feature} & \textbf{DRAM} & \textbf{SRAM} \\
      \midrule
      Full form & Dynamic RAM & Static RAM \\
      Refresh & Needs periodic refresh & No refresh \\
      Speed & Slower & Faster \\
      Density / cost & Denser, cheaper & Less dense, costlier \\
      Typical use & Main memory & Cache memory \\
      \bottomrule
    \end{tabular}
  \end{center}
  \vspace{0.3cm}
  \muted{DRAM: many cheap cells. SRAM: fewer, faster cells.}
\end{frame}

\begin{frame}{ROM --- Non-Volatile and Read-Only}
  \begin{itemize}
    \item \key{ROM} is \key{Read-Only Memory}.
    \item \key{Non-volatile}: it keeps its contents without power.
    \item \key{Read-only in normal operation}; writing needs a special
      programming process.
    \item It commonly stores \key{firmware} such as the start-up code.
    \item RAM and ROM are both primary memory; they differ in volatility and
      writeability.
  \end{itemize}
  \begin{alertblock}{Trap (TRAP-02)}
    \bad{ROM is not freely read and write.} Normally you can only read it.
  \end{alertblock}
\end{frame}

\begin{frame}{The ROM Family}
  \begin{center}
    \begin{tabular}{ll}
      \toprule
      \textbf{Type} & \textbf{Key property} \\
      \midrule
      PROM & Programmable once; cannot be erased \\
      EPROM & Erased by ultraviolet (UV) light \\
      EEPROM & Erased and rewritten electrically \\
      Flash & EEPROM-derived; erased in blocks \\
      \bottomrule
    \end{tabular}
  \end{center}
  \vspace{0.3cm}
  \muted{PROM once \quad EPROM by UV \quad EEPROM electrically \quad Flash in blocks.}
\end{frame}

\begin{frame}{Cache --- Fast Memory Near the CPU}
  \begin{itemize}
    \item \key{Cache} is a small, very fast memory between the CPU and main
      memory.
    \item It holds frequently or recently used data and instructions.
    \item Purpose: reduce the CPU's wait for slower main memory.
    \item Cache is \key{faster than main memory} (PYQ-03).
    \item Cache is working memory, \bad{not} a backup or permanent store.
  \end{itemize}
  \begin{block}{Exam cue}
    Speed order: cache $>$ main memory $>$ secondary storage.
  \end{block}
\end{frame}

\begin{frame}{Cache Levels: L1, L2, L3}
  \begin{itemize}
    \item \key{L1}: smallest and fastest, closest to the core.
    \item \key{L2}: larger and a little slower than L1.
    \item \key{L3}: largest and slowest of the three, typically shared
      across cores.
    \item All cache levels are faster than main memory.
  \end{itemize}
  \begin{block}{Exam cue}
    Size and distance: \key{L1\,$<$\,L2\,$<$\,L3}. Speed: reverse.
  \end{block}
\end{frame}

\begin{frame}{Cache Hit, Miss, and Hit Ratio}
  \begin{itemize}
    \item \key{Hit}: the CPU finds the required data in the cache.
    \item \key{Miss}: the data is not in the cache, so it is fetched from
      main memory.
    \item \key{Hit ratio} $=$ hits $\div$ total accesses.
    \item Example: 100 accesses with 90 hits give a hit ratio of
      $90/100 = 0.90$.
    \item A higher hit ratio means less time waiting for main memory.
  \end{itemize}
\end{frame}

\begin{frame}{RAM vs ROM vs Cache}
  \begin{center}
    \begin{tabular}{lll}
      \toprule
      \textbf{Aspect} & \textbf{RAM} & \textbf{ROM} \\
      \midrule
      Volatility & Volatile (TRAP-01) & Non-volatile \\
      Normal access & Read and write & Read-only (TRAP-02) \\
      Typical role & Main memory & Firmware / start-up code \\
      Speed & Faster than secondary & Slower than cache \\
      \bottomrule
    \end{tabular}
  \end{center}
  \vspace{0.2cm}
  \muted{Cache sits between the CPU and RAM and is the fastest of the three.}
\end{frame}

\begin{frame}{Common Traps in This Topic}
  \begin{itemize}
    \item \bad{TRAP-01}: RAM is volatile; ROM is non-volatile.
    \item \bad{TRAP-02}: ROM is read-only in normal operation.
    \item Cache is \key{faster} than main memory, not slower (PYQ-03).
    \item \key{SRAM} goes to cache; \key{DRAM} goes to main memory.
    \item Cache is \bad{not} a backup --- it is temporary working memory.
  \end{itemize}
\end{frame}

\begin{frame}{Lesson 08: Recall}
  \begin{itemize}
    \item RAM is volatile; ROM is non-volatile and normally read-only.
    \item DRAM: denser/cheaper, needs refresh, main memory; SRAM: faster,
      no refresh, cache.
    \item ROM family: PROM (once), EPROM (UV), EEPROM (electrical), flash
      (block erase).
    \item L1 smallest/fastest, L3 largest/shared; all faster than RAM.
    \item Hit ratio $=$ hits $\div$ total accesses.
  \end{itemize}
\end{frame}
\end{document}
